\documentclass[12pt]{article}
\usepackage{../../../assignment}
\usepackage{../../../math}
\usepackage{../../../uark_colors}
\hypersetup{
  colorlinks = true,
  allcolors = ozark_mountains,
  breaklinks = true,
  bookmarksopen = true
}

\newcommand{\answer}[1]{{\color{blue_winged_teal}\textbf{Answer:} #1}}
\newcommand{\pts}[1]{{\color{zinc500}(#1 pts)}}

\begin{document}
\begin{center}
  {\Huge\bf Final - Fall 2025}

  \smallskip
  {\large\it ECON 5783 — University of Arkansas}
\end{center}
\medskip

\begin{enumerate}
  \item A researcher wants to evaluate whether a school having an arts program benefits children's educational outcomes. 
  The researcher surveys schools and records their average GPA ($Y_i$) and whether or not they have an arts program ($D_i$).
  \begin{enumerate}
    \item \pts{10} In this context, describe in words what the potential outcomes, $Y_i(0)$ and $Y_i(1)$, are.
    
    \answer{
      $Y_i(0)$ is the counterfactual outcome if a particular school has no arts program and $Y_i(1)$ is the counterfactual outcome if that school had an arts program.
    }
    
    \medskip
    \item \pts{10} Do you think the difference-in-means estimator would be biased in this context? 
    If so, which direction do you the treatment effect estimate is biased?

    \answer{Typically, we think the presence of arts programs are correlated with school resources and average family incomes. Since both of these variables have a positive effect on the educational outcomes, there is a positive bias in our estimate.}
    
    \medskip
    \item \pts{10} Why might a table of summary statistics consisting of sample means of observables, broken down by $D_i = 0$ and $D_i = 1$, be useful in this setting?
    
    \answer{
      Demonstrate which variables are correlated with treatment (having an arts program). 
      This table can tell us what we might need to control for.  
    }
  \end{enumerate}
\end{enumerate}

\bigskip\bigskip
%% https://www.christopher-neilson.com/work/documents/UPK/w33038.pdf
\noindent The following questions are about Rosenthal and Strange (2008, Journal of Urban Economics). 
They want to understand if having a higher density of firms in an area increases the average wages of workers. 

\begin{enumerate}
  \setcounter{enumi}{1}
  \item \pts{10} Consider regressing the average wages of a worker in a city on the density of firms in the area. 
  Say you find a large, positive estimate.
  Using the concept of omitted variables bias, give an example of why we should not believe we identified the causal effect of density.

  \answer{
    There might be some factors in a city that are conducive for firm's productivity (e.g. presence of a high-quality college). 
    This would lead to a large density of firms sorting there as well as higher worker wages.
  }

  \medskip
  \item \pts{10} Rosenthal and Strange use as an instrument the proportion of soil in the city that contains sedimentary rock. 
  The idea being that sedimentary soil is more difficult to build on which causes less density. 
  Argue why the exclusion restriction might be reasonable in this context.

  \answer{
    It is likely that sedimentary soil does not have any impact on firm productivity outside of by limiting the density of buildings. 
    Therefore, it is unlikely that the instrument is correlated with the reduced form's residual.
  }

  \medskip
  \item \pts{10} The first-stage regression of density on the sedimentary rock share has a $F$-statistic of 1.907. 
  Describe why this might cause a problem when conducting inference on our 2SLS estimate.

  \answer{
    Since the 2SLS estimate is the ratio of the reduced form over the first-stage, a small first-stage means dividing by a number close to 0. 
    Hence this makes our estimate very sensitive to noise in the first-stage estimate.
  }
\end{enumerate}



\newpage
\noindent The following questions are about Leah Boustan's work on court-enforced school desegregation (2012, AEJ Applied).
They ask whether forced school desegregation in the city-center led to a larger decline in the housing price in the city-center.
The unit of observation is a metro area and they are observed in 1960, 1970, and 1980 (post-period).
Below is the primary DID estimates.

\bigskip
\begin{enumerate}
  \setcounter{enumi}{4}
  \item \pts{10} The row with $\Delta$ 1970--1980 presents the difference-in-differences estimate. 
  What do we have to assume about the cities that were and were not placed under court orders in order to belive the 5.8\% decrease in home prices estimate?
  
  \answer{
    We need to assume the parallel counterfactual trends assumption.
    Namely, we have to believe that the trend in housing prices in the treated cities had they not been treated would be the same, on average, as the non-treated cities.
  }

  \medskip
  \item \pts{10} Why is the row labeled ``$\Delta$ 1960--1970'' helpful for justifying the causal assumption?
  
  \answer{
    Prior to the court-ordered desegregation orders, it seems that the treated and untreated had a similar change in home prices in the city-center.
    This helps make us more comfortable that the counterfactual trend from 1970 and 1980 would be similar for these cities.
  }

  \medskip
  \item \pts{10} Boustan suggests that cities that had larger shares of Black residents might have had larger suburbanization rates during this time (`White flight'). 
  Moreover, cities that put under court-ordered forced desegregation had larger 
  Why might this represent concerns for our parallel trends assumption?

  \answer{
    In the absence of school desegregatin, treated cities might have had a higher (counterfactual) price decline in the city-center because of higher rates of White flight. 
    This would create a violation of the (unconditional) parllel trends assumption.
  }

  \medskip
  \item \pts{10} Say you are willing to assume parallel trends \emph{conditional on} the share of Black residents in 1970, i.e. $\Delta Y_i(0) \Perp D_i \ \vert \ X_i = x$.
  Describe how you would estimate a treatment effect using the regression adjustment strategy. 

  \answer{
    \begin{enumerate}
      \item Regress $\Delta Y_i$ on (functions of) $X_i$ using the $D_i = 0$ units to fit a model for $\Delta Y_i(0)$.
      
      \item Impute $\widehat{\Delta Y}_i(0)$ for $D_i = 1$ units
      
      \item Take average of $\Delta Y_i - \widehat{\Delta Y}_i(0)$ for treated units to estiamate the ATT.
    \end{enumerate}
  }
\end{enumerate}

\bigskip
\begin{table}[hb]
\centering
\caption{School Desegregation and Relative City Housing Prices at the District Border, 1960--1980}
\label{tab:desegregation}
\begin{tabular}{lccc}
\toprule
Dependent variable & Placed under & Not placed under & \\
$= \ln(\text{housing value})$; & court order & court order & \\
coefficient on & during 1970s & during 1970s & Difference \\
\midrule
% 1970 & $-0.047***$ & $-0.026*$ & $-0.021$ \\
%      & (0.014) & (0.015) & (0.020) \\[0.5em]
% 1980 & $-0.097***$ & $-0.023$ & $-0.073**$ \\
%      & (0.028) & (0.022) & (0.035) \\[0.5em]
$\Delta$ 1970--1980 & $-0.065**$ & $-0.007$ & $-0.058**$ \\
     & (0.024) & (0.015) & (0.028) \\[0.5em]
\textit{Pre-trend:} & & & \\
$\Delta$ 1960--1970 & $-0.023*$ & $-0.022$ & $-0.001$ \\
     & (0.013) & (0.017) & (0.022) \\
\bottomrule
\end{tabular}
\end{table}

\end{document}
